% STATUS: VERIFIED / AUDITED
% DEPENDS ON: P02, P03, P04, P09 VERIFIED / AUDITED
% INHERITED FOUNDATION: P01
% MODEL: canonical DM
% FOLD REGIME: a > k+h
% BIFURCATION PARAMETER: g
% ROLE: full-system positive saddle-node and stable-coexistence theorem
% NUMERICAL DEPENDENCIES: none
% NEXT DEPENDENCY: P11
% BASIN SCOPE: synchronous threshold + local stable manifold only unless
%              stronger statements are explicitly proved
% VERSION: P10 v1.1; 2026-10-03; metadata-only A10 audit closure
% CLAIMS: C070, C071; C147 universal claim refuted by an exact witness
% LOCAL BASIN REFINEMENT: two basin sides in a sufficiently small saddle
% neighborhood are proved below; no complete global basin boundary claimed.
% FROZEN ROLE: supporting proposition/lemma S4; no scope amendment needed.
% EXCLUSIONS: no P05/P07/P08 theorem input; no P06 or P11--P18 theorem.
% AUDIT FILE: audits/A10_saddle_node_and_bistability.tex
% AUDIT VERDICT: VERIFIED; required corrections: NONE
% LAYER VERDICTS: fold geometry/spectrum VERIFIED;
% full-system saddle-node VERIFIED; branch stability/bistability VERIFIED;
% coexistence set/counterexample VERIFIED; basin/manifold VERIFIED;
% scope/dependency VERIFIED. OVERALL: VERIFIED.
% AUDITED VERSION: P10 v1.0; 2026-10-03; 32560 bytes; 707 lines
% AUDITED FULL-FILE HASH:
%   a28ea2b4ca96f381f06fa82849d4289fcc0e9c41e646ca66595621922f385a07
% AUDITED MATHEMATICAL-BODY HASH:
%   5cfd8afb91b7f23575e1b4bd9c7c28f3b433affefae9bb0d5f7700d030fcba81
% CLOSURE DATE: 2026-10-03
% CURRENT VERSION: P10 v1.1; metadata only; no mathematical revision
% CURRENT FULL-FILE HASH: recorded externally in project_manifest.md and
% the closure report, avoiding a self-referential full-file hash.
% CURRENT MATHEMATICAL-BODY HASH:
%   5cfd8afb91b7f23575e1b4bd9c7c28f3b433affefae9bb0d5f7700d030fcba81
% BODY SCOPE: first line-anchored section command through EOF, all bytes.
% PRESERVATION: mathematical body is byte-identical to exact audited v1.0.
% STATUS NOTE: historical DRAFT/audit-pending body text is preserved;
% this header and the CURRENT manifest closure govern lifecycle status.
% VALIDATION PROVENANCE: A10 authenticated/inventoried the original developer
% validation package but did NOT rerun it or claim reproduction of its results.
% A10 analytical reconstruction supplies verification; its independent finite
% checker is separate corroboration. No checker is rerun for closure.
% BASIN SCOPE PRESERVED: synchronous-ray threshold; local full-space disk
% boundary only relative to its neighborhood V; no complete global partition.
% P11 GATE: cleared to begin after this closure; not begun in this task.

\section{P10: Saddle-node bifurcation and true bistable region}
\label{sec:P10}

\noindent\textbf{Status: DRAFT; independent A10 audit pending.}
Fix $a,k,c,h>0$, $n\ge5$, $1\le q<n$, and vary $g>0$ in the
canonical DM field \eqref{eq:canonical-dm}. Put
\begin{equation}
 A=k+h,\quad b=a-A>0,\quad g_c=a+c,\quad
 g_f=g_c+\frac{b^2}{4A},\quad r_f=\frac b{2A},\quad s_f=1+r_f.
 \label{eq:P10-data}
\end{equation}
Here $r$ is the synchronous amplitude and $b$ is a local abbreviation.
\emph{Notation reconciliation:} P02's roots $s_-,s_+$ are the amplitudes
called $r_-,r_+$ here. P02's repeated-root amplitude $s_f$ in
\eqref{eq:P02-gf} is our $r_f$; P10's $s_f$ denotes its denominator
$1+r_f$, following P03/P04's denominator convention. No earlier notation
or file is changed. All neighborhoods below may depend on the fixed
parameters and on $(n,q)$; no uniform spectral gap as $n\to\infty$ is claimed.

The direct inputs are P02's invariant restriction and root formulas,
P03's original-field Jacobian and origin stability, P04's complete
Fourier spectrum, and P09's origin branch organization. P01 supplies
smoothness, uniqueness and global nonnegative existence. We also use
the real implicit-function theorem, the usual smooth local stable/unstable
manifold theorem, and linearized stability/instability. These standard
local results introduce no new project-module dependency or literature work.

\subsection{Exact fold geometry}
\label{slot:P10-1}

P02 gives, for $r>-1$,
\begin{equation}
 F(r\one,g)=f(r,g)\one,\qquad
 f(r,g)=ar-Ar^2+\frac{(c-g)r}{1+r}
       =\frac r{1+r}(\mu+br-Ar^2),\quad \mu=g_c-g.
 \label{eq:P10-scalar}
\end{equation}
The nonzero equilibrium equation is $Ar^2-br-\mu=0$.
Completing the square, or using its discriminant, verifies independently
\begin{equation}
 g_f=c+\frac{(a+A)^2}{4A},\qquad
 s_f=1+\frac{a-A}{2A}=\frac{a+A}{2A},\qquad
 g_f-g_c=\frac{b^2}{4A}>0.
 \label{eq:P10-fold-identities}
\end{equation}
Define the state and parameter coordinates
\begin{equation}
 u=r-r_f,\qquad \nu=g_f-g.
 \label{eq:P10-coordinates}
\end{equation}
Then $\mu=\nu-b^2/(4A)$, $br_f-Ar_f^2=b^2/(4A)$, and
$b-2Ar_f=0$. Expanding in $u$ therefore proves, with no approximation,
\begin{equation}
 \boxed{\mu+br-Ar^2=\nu-Au^2,\qquad
 \dot u=\frac{r_f+u}{s_f+u}(\nu-Au^2).}
 \label{eq:P10-exact-fold}
\end{equation}
Choose $|u|<r_f/2$. Both amplitude and denominator are positive.
For $\nu<0$ there is no nonzero synchronous equilibrium anywhere on
$r>-1$; in particular none near $r_f$. For $\nu=0$ there is one
double positive scalar equilibrium $r_f$. For small $\nu>0$ there
are two distinct nearby positive roots
\begin{equation}
 \boxed{r_\pm=r_f\pm\sqrt{\nu/A}=s_\pm\ \text{of P02}.}
 \label{eq:P10-local-roots}
\end{equation}
Both are positive exactly when $0<\nu<b^2/(4A)$ within this real
root regime. The origin always exists, but is excluded from the
chosen fold state neighborhood.

\subsection{Complete fold spectrum and the real transverse reduction}

For a synchronous state let $s=1+r$ and import P03/P04:
\[
 d=a-2kr-\frac{cr}{s^2}-\frac gs-hr,\qquad
 C=\frac cs,\quad G=\frac{gr}{s^2},\quad H=hr.
\]
\begin{align}
 J(r,g)&=dI+CP+GP^{-q}-HP^q,\notag\\
 \lambda_m&=d+Ce^{-\ii\theta_m}+Ge^{\ii q\theta_m}-He^{-\ii q\theta_m},
 \label{eq:P10-spectrum}
\end{align}
where $\theta_m=2\pi m/n$ and $0\le m<n$.
Colliding offsets are added, as in P03; none is removed.
At a positive equilibrium P04 gives
\begin{equation}
 \lambda_0=f_r(r,g)=\frac r{1+r}(b-2Ar).
 \label{eq:P10-uniform}
\end{equation}
Thus $\lambda_0^f=0$, also obtained by differentiating
\eqref{eq:P10-exact-fold} at $(u,\nu)=(0,0)$.

\begin{lemma}[Automatic transverse stability at every positive fold]
\label{lem:P10-gap}
At the fold, with subscript $f$ denoting evaluation there,
\begin{equation}
 \boxed{G_f-H_f=r_f\left(k+\frac{4A^2c}{(a+A)^2}\right)>0.}
 \label{eq:P10-Lf}
\end{equation}
The exact complex fold spectrum is \eqref{eq:P10-spectrum}, with
$d_f=-C_f-G_f+H_f$. In particular,
\begin{align}
 \Ree\lambda_m^f
 &=C_f(\cos\theta_m-1)+(G_f-H_f)(\cos(q\theta_m)-1),
 \label{eq:P10-fold-real}\\
 \Imm\lambda_m^f&=-C_f\sin\theta_m+(G_f+H_f)\sin(q\theta_m).
 \label{eq:P10-fold-imag}
\end{align}
For every $m\ne0$,
\begin{equation}
 \boxed{\Ree\lambda_m^f\le-\delta_f<0,\qquad
 \delta_f=\frac c{s_f}\left(1-\cos\frac{2\pi}{n}\right)>0.}
 \label{eq:P10-gap}
\end{equation}
The only eigenvalue on the imaginary axis is zero. It has algebraic
and geometric multiplicity one and real eigenspace
$\operatorname{span}\{\one\}$.
\end{lemma}
\begin{proof}
Using \eqref{eq:P10-fold-identities}, directly
\begin{equation}
 \frac{g_f}{s_f^2}-h
 =\frac{c+A s_f^2}{s_f^2}-h
 =k+\frac c{s_f^2}=k+\frac{4A^2c}{(a+A)^2}>0.
 \label{eq:P10-effective}
\end{equation}
Multiplication by $r_f>0$ proves \eqref{eq:P10-Lf}.
Since $0=\lambda_0^f=d_f+C_f+G_f-H_f$, taking real and imaginary
parts of the complete P04 formula proves the two identities.
For $1\le m<n$, set $j=\min(m,n-m)$. The angle
$2\pi j/n$ belongs to $[2\pi/n,\pi]$, so
$\cos\theta_m\le\cos(2\pi/n)<1$. The first cosine difference
is strictly negative, while the second is nonpositive. This proves
the margin. For even $n$ the real singleton $m=n/2$ is included;
for odd $n$ the same argument ends before $\pi$. Collisions
$q=1,n-1,n/2$ and $\gcd(n,q)>1$ change neither argument.
Even if $q\theta_m$ is a multiple of $2\pi$, the first term is
strictly negative. P04's unitary diagonalization lists exactly $n$
eigenvalues with algebraic multiplicity. Only label zero can vanish,
so its multiplicities are one and its eigenvector is $\one$.
No simplicity of the other numerical eigenvalues is asserted.
\end{proof}

\begin{proposition}[Full-system local equilibrium completeness]
\label{prop:P10-reduction}
In a state neighborhood of $r_f\one$ and parameter interval about
$g_f$, every full-system equilibrium is synchronous. The entire local
equilibrium equation reduces exactly to $\nu-Au^2=0$.
\end{proposition}
\begin{proof}
Let $v=\one$, $w^{\mathsf T}=n^{-1}\one^{\mathsf T}$,
$E=\{y:w^{\mathsf T}y=0\}$, and $\Pi_\perp=I-vw^{\mathsf T}$.
Write uniquely $x=r\one+y$, $y\in E$. The circulant $J_f$
preserves $E$, and its restriction to the complexification of $E$
has exactly the nonuniform eigenvalues. Lemma~\ref{lem:P10-gap}
therefore makes $J_f|_E:E\to E$ an invertible real map.

For the smooth $E$-valued map
$T(r,y,g)=\Pi_\perp F(r\one+y,g)$ one has
$D_yT(r_f,0,g_f)=J_f|_E$. The real implicit-function theorem gives
a product neighborhood and a unique smooth transverse solution
$y=\psi(r,g)$. Invariance of synchrony gives $T(r,0,g)=0$ for
every nearby $r,g$. After shrinking the product, uniqueness forces
$\psi\equiv0$. Every full equilibrium solves $T=0$ and hence has
$y=0$. Conversely, at $y=0$ its scalar equation is $f=0$.
Here $r>0$ and $1+r>0$, so \eqref{eq:P10-exact-fold} makes it
precisely $\nu-Au^2=0$. A smaller state ball times a parameter
interval lies in the IFT product. This is a full equilibrium reduction,
not an inference from the presence of synchronous solutions alone.
\end{proof}

\subsection{Full-system nondegeneracy and exact center dynamics}

\begin{proposition}[Original-vector-field saddle-node coefficients]
\label{prop:P10-coefficients}
With the above $v,w$, one has $J_fv=0$, $w^{\mathsf T}J_f=0$,
$w^{\mathsf T}v=1$, and
\begin{align}
 \alpha_{\rm SN}:=w^{\mathsf T}F_g(r_f\one,g_f)
 &=-\frac{r_f}{s_f}<0,\label{eq:P10-alpha}\\
 w^{\mathsf T}F_{xx}(r_f\one,g_f)[v,v]
 &=-\frac{2Ar_f}{s_f}\ne0,\label{eq:P10-contraction}\\
 \beta_{\rm SN}:=\tfrac12w^{\mathsf T}F_{xx}[v,v]
 &=-\frac{Ar_f}{s_f}<0.\label{eq:P10-beta}
\end{align}
Derivatives hold all other physical coefficients fixed. With $\nu$
as parameter the parameter coefficient has the opposite sign,
$+r_f/s_f$.
\end{proposition}
\begin{proof}
The row and column sums of $J_f$ are both $\lambda_0^f=0$, giving
the right and left nullvectors and their normalization. In the original
vector field, $(F_g)_i=-x_i/(1+x_{i+q})$, hence its average at the
fold is $-r_f/s_f$.

Independently differentiate each original summand twice along
$x=(r_f+t)\one$ at fixed $g=g_f$. The five componentwise second
derivatives, retaining all terms even at coincident offsets, are
\[
 0,\qquad -2k,\qquad -\frac{2c}{s_f^3},\qquad
 \frac{2g_f}{s_f^3},\qquad -2h.
\]
Their sum is $-2A+2(g_f-c)/s_f^3$. Since $g_f-c=A s_f^2$,
this equals $-2A+2A/s_f=-2Ar_f/s_f$ in every component.
Averaging produces no extra factor of $n$. The Taylor coefficient
of $u^2$ is half this full contraction; the derivatives themselves
contain no Taylor factorial.
\end{proof}

\begin{theorem}[Nondegenerate full-system positive saddle-node]
\label{thm:P10-saddle-node}
For every fixed canonical $a>k+h$, $c,k,h>0$ and admissible $(n,q)$,
the equilibrium $r_f\one$ at
\[
 \boxed{g=g_f=a+c+\frac{(a-k-h)^2}{4(k+h)}}
\]
is a nondegenerate full-system saddle-node. It has a simple uniform
zero eigenvalue and strictly stable transverse spectrum. Locally there
are no equilibria for $g>g_f$, one double equilibrium for $g=g_f$,
and precisely the two positive equilibria \eqref{eq:P10-local-roots}
for $g<g_f$ sufficiently close. The exact center restriction is
orbitally equivalent, by a positive change of time, to
\begin{equation}
 \boxed{\frac{du}{d\tau}=\nu-Au^2,\qquad \frac{d\nu}{d\tau}=0.}
 \label{eq:P10-normal-form}
\end{equation}
\end{theorem}
\begin{proof}
Smoothness holds in an open neighborhood contained in $\PosOrth$
and with $g>0$. Lemma~\ref{lem:P10-gap} supplies the simple kernel
and stable complement. Proposition~\ref{prop:P10-coefficients}
verifies both nonzero coefficients of the standard saddle-node
nondegeneracy test. Proposition~\ref{prop:P10-reduction} proves
local equilibrium completeness, so no nonsynchronous branch is omitted.

For completeness the normal form is exact on a known center manifold.
At fixed $g_f$ the synchronous line is invariant and tangent to the
entire one-dimensional center eigenspace. For the parameter extension
$(\dot x,\dot\nu)=(F(x,g_f-\nu),0)$, the linearization at the fold,
in coordinates $(u,y,\nu)$, is
\begin{equation}
 \dot u=(r_f/s_f)\nu,\qquad \dot y=(J_f|_E)y,\qquad \dot\nu=0.
 \label{eq:P10-extended-linear}
\end{equation}
Its center block is a nilpotent two-by-two block: unlike P09 at the
origin, the parameter derivative here is nonzero. The generalized
center subspace is precisely $y=0$. Synchrony times the parameter
interval is smooth, invariant and tangent to this subspace, hence an
exact local two-dimensional center manifold. Its exact restriction
is \eqref{eq:P10-exact-fold} with $\dot\nu=0$.
For $|u|<r_f/2$ the time factor satisfies
\[
 \frac{r_f+u}{s_f+u}\ge\frac{r_f/2}{1+3r_f/2}>0.
\]
Consequently $d\tau/dt=(r_f+u)/(s_f+u)$ proves
\eqref{eq:P10-normal-form}. This is orbital equivalence on the
center restriction, not a global decoupling of the full field.
The quadratic equation gives the stated branches and sides.
At the fold itself the center flow is negative for $u\ne0$;
it attracts from $u>0$ and escapes from $u<0$, so the fold is
not a locally asymptotically stable equilibrium in the full space.
\end{proof}

\subsection{Both branches and all Fourier modes}
\label{slot:P10-2}

\begin{proposition}[Exact uniform signs throughout the two-root interval]
\label{prop:P10-branch-signs}
For every $g_c<g<g_f$, define $D$ as in P02 and $\Delta=\sqrt D$.
Then
\begin{align}
 0<\Delta&=\sqrt{b^2+4A(a+c-g)}=2\sqrt{A(g_f-g)}<b,\notag\\
 r_\pm&=\frac{b\pm\Delta}{2A},\qquad
 b-2Ar_\pm=\mp\Delta,\label{eq:P10-discriminant}\\
 \boxed{\lambda_0(r_\pm)=\mp\frac{r_\pm}{1+r_\pm}\Delta.}
 \label{eq:P10-branch-eigenvalues}
\end{align}
Thus the lower branch is always unstable in the uniform direction,
whereas the upper uniform eigenvalue is always strictly negative.
Any upper-branch loss of stability within this open interval must
involve a nonuniform mode.
\end{proposition}
\begin{proof}
P02's discriminant factorization is also obtained directly by inserting
$g_f=g_c+b^2/(4A)$. The strict parameter inequalities give $0<\Delta<b$,
so both roots are positive. Substitution into $b-2Ar_\pm$ cancels
$b$ and leaves $\mp\Delta$. Equation~\eqref{eq:P10-uniform}
gives the exact signs, and the positive lower eigenvalue implies
nonlinear instability. This last implication needs no hyperbolicity
of the other lower-branch eigenvalues.
\end{proof}

\begin{corollary}[Near-fold full branch stability]
\label{cor:P10-near-stability}
There exists $\varepsilon_{\rm sp}>0$ such that both branches for
$0\le\nu<\varepsilon_{\rm sp}$ satisfy
\[
 \Ree\lambda_m(r_\pm,g_f-\nu)<-\delta_f/2\quad(m\ne0).
\]
For $0<\nu<\min(\varepsilon_{\rm sp},b^2/(4A))$, the upper
branch is Hurwitz and locally exponentially stable. The lower branch
is a hyperbolic saddle with exactly one simple positive real eigenvalue
and $n-1$ stable real dimensions. Its local stable manifold has dimension
$n-1$ and codimension one, and its unstable manifold has dimension one.
\end{corollary}
\begin{proof}
Every labeled eigenvalue in \eqref{eq:P10-spectrum} is continuous
in $(r,g)$ near $(r_f,g_f)$. There are finitely many labels.
A common neighborhood makes each transverse real part differ from
its fold value by less than $\delta_f/2$. Both continuous roots
$r_\pm(\nu)$ tend to $r_f$ as $\nu\downarrow0$; smoothness in
$\nu$ at zero is not needed. Shrink $\varepsilon_{\rm sp}$ so both
branches lie in that neighborhood and remain positive. Combine the
margin with \eqref{eq:P10-branch-eigenvalues}. Only label zero has
positive real part on the lower branch, so its eigenvalue is simple.
Real conjugate-pair blocks and real singleton modes together give
$n-1$ stable dimensions, even if stable numerical eigenvalues repeat.
P04's nonlinear stability implications and the smooth local
stable/unstable manifold theorem apply.
\end{proof}

\subsection{The nonempty open bistability set}
\label{slot:P10-3}

\begin{theorem}[Stable coexistence and openness in physical parameters]
\label{thm:P10-bistability}
For each fixed $a,k,c,h>0$ with $a>k+h$ and each admissible $(n,q)$,
there is $\varepsilon>0$ with
\begin{equation}
 (g_f-\varepsilon,g_f)\subset(g_c,g_f)
 \label{eq:P10-bistable-interval}
\end{equation}
such that throughout this interval the full system has the locally
exponentially stable origin, a positive index-one saddle $r_-\one$,
and the locally exponentially stable positive equilibrium $r_+\one$.
At every strict point of a sufficiently small such interval this
configuration persists on an open neighborhood in $(0,\infty)^5$
with coordinates $(a,k,c,g,h)$ and fixed $(n,q)$.
Here ``bistability'' means coexistence of at least two locally stable
equilibria; it does not mean that these are the only attractors.
\end{theorem}
\begin{proof}
At the origin P03 gives
\[
 \Ree\lambda_m(0)=g_c-g-c(1-\cos\theta_m)\le g_c-g<0
 \quad(g_c<g<g_f).
\]
Thus the origin is Hurwitz and locally exponentially stable in
$\Domain$, and also relative to $\Orth$ using P01 invariance.
Take $0<\varepsilon<\min(\varepsilon_{\rm sp},b^2/(4A))$ and
apply Corollary~\ref{cor:P10-near-stability}. The interval is nonempty
and lies strictly above $g_c$.

At a strict witness all three Jacobians are invertible and hyperbolic.
The implicit-function theorem continues each equilibrium smoothly
under changes of all five physical coefficients. The origin stays
identically zero. Alternatively the positive roots are explicitly
smooth because $D>0$; uniqueness identifies them with the IFT
continuations. Their separation and strict positivity persist after
shrinking the parameter neighborhood. Spectral continuity preserves
both Hurwitz properties and the lower saddle's unstable dimension.
The inequalities $a>A$, $g_c<g<g_f$ and positivity of every coefficient
also persist. Such neighborhoods can depend on $(n,q)$; no common
all-dimension persistence neighborhood for the index-one configuration
is asserted.
\end{proof}

\begin{proposition}[Exact coexistence set and a simple sufficient criterion]
\label{prop:P10-bistability-set}
For fixed $(n,q)$ define the open two-positive-root domain
\[
 \mathcal D=\{\param\in(0,\infty)^5:a>k+h,\ g_c<g<g_f\}.
\]
For $\param\in\mathcal D$, evaluate \eqref{eq:P10-spectrum} at its
upper root and define
\begin{equation}
 \boxed{\mathcal B_{n,q}=
 \left\{\param\in\mathcal D:
 \max_{0\le m<n}\Ree\lambda_m(r_+,g)<0\right\}.}
 \label{eq:P10-B}
\end{equation}
This set is open and nonempty. It is exactly the set in $\mathcal D$
where the origin and upper synchronous equilibrium are both locally
exponentially stable. Every point has an unstable lower positive
branch, although its unstable dimension is asserted to be one only
in the near-fold neighborhood above. A sufficient condition for
membership in $\mathcal B_{n,q}$ is
\begin{equation}
 \boxed{G\ge H\quad\Longleftrightarrow\quad
 \frac{g}{(1+r_+)^2}\ge h.}
 \label{eq:P10-sufficient}
\end{equation}
This condition allows equality.
\end{proposition}
\begin{proof}
The positive upper root and all its labeled eigenvalues are continuous
(and smooth before taking the finite maximum) on $\mathcal D$.
The strict finite set of inequalities defines an open subset of
$(0,\infty)^5$, hence of $\R^5$. Nonemptiness follows from
Theorem~\ref{thm:P10-bistability}. The origin is automatically Hurwitz
on $\mathcal D$. A Hurwitz Jacobian gives nonlinear local exponential
stability. Conversely, a smooth equilibrium satisfying a local
exponential estimate has a Jacobian semigroup satisfying that estimate
by differentiation of the flow at the equilibrium, so all eigenvalues
have strictly negative real part. Thus the definition is exact for
exponential stable coexistence, not a classification of possible
nonhyperbolic asymptotically stable boundary points or other attractors.

Subtracting the uniform row sum in \eqref{eq:P10-spectrum} gives,
at either branch and in fact at any synchronous state,
\begin{equation}
 \Ree\lambda_m=\lambda_0+C(\cos\theta_m-1)
                         +(G-H)(\cos(q\theta_m)-1).
 \label{eq:P10-difference}
\end{equation}
On the upper branch $\lambda_0<0$ and $C>0$. If $G-H\ge0$,
every mode has negative real part, including $m=0$; for $m\ne0$
the $C$ contribution is strictly negative. Since $r_+>0$, dividing
$G-H$ by it gives the physical-coefficient condition. The lower
uniform sign follows from Proposition~\ref{prop:P10-branch-signs}.
\end{proof}

One can give an explicit, dimension-independent sufficient interval
for the two stable equilibria (not necessarily for the saddle's index).
Set
\begin{equation}
 M=g_f-hs_f^2=c+ks_f^2>0,\qquad
 t_*:=\min\left\{\frac{r_f}{2},\frac{M}{4hs_f},
                         \sqrt{\frac{M}{2(A+h)}}\right\}>0.
 \label{eq:P10-explicit-margin}
\end{equation}
For $0<\nu<A t_*^2$, put $t=\sqrt{\nu/A}<t_*$. Then
$g=g_f-At^2$, $r_+=r_f+t$, and directly
\begin{equation}
 g-h(1+r_+)^2=M-2hs_ft-(A+h)t^2>0.
 \label{eq:P10-explicit-condition}
\end{equation}
Indeed each of the last two terms is strictly less than $M/2$.
Also $At_*^2\le b^2/(16A)<g_f-g_c$, so this whole interval lies
in $\mathcal D$ and in $\mathcal B_{n,q}$ for every admissible
$(n,q)$. This supplies a quantitative alternative to the continuity
argument for stable coexistence. Equation~\eqref{eq:P10-effective}
also proves directly that $G>H$ persists on the upper branch near
the fold. Intersecting this interval with that of
Corollary~\ref{cor:P10-near-stability} retains the index-one conclusion.

\begin{remark}[An exact obstruction to whole-interval stability]
\label{rem:P10-counterexample}
Take
\begin{equation}
 n=6,\quad q=1,\qquad
 (a,k,c,g,h)=\left(6,\frac1{10},\frac1{10},\frac{101}{10},\frac9{10}\right).
 \label{eq:P10-counterexample}
\end{equation}
Here $A=1$, $b=5$, $g_c=61/10$, $g_f=247/20$, $\mu=-4$,
and the equilibrium polynomial is $(r-1)(r-4)$.
At $r_+=4$ the original-field equilibrium residual is exactly zero,
and
\[
 d=-\frac{109}{250},\quad C=\frac1{50},\quad
 G=\frac{202}{125},\quad H=\frac{18}{5},\quad
 G-H=-\frac{248}{125}.
\]
Consequently $\lambda_0=-12/5<0$, but the real mode $m=3$ has
\begin{equation}
 \lambda_3=d-C-G+H=\frac{191}{125}>0.
 \label{eq:P10-counterexample-mode}
\end{equation}
The origin is stable and the upper equilibrium unstable at this
strict interior point of the two-root interval. This refutes universal
upper-branch stability, hence universal origin/upper stable coexistence
on that interval. At the lower root $r_-=1$, one has
$\lambda_0=3/2$ and $\Ree\lambda_1=\Ree\lambda_5=53/80>0$;
the other three real parts are negative. Thus its unstable dimension
is three here, also precluding a whole-interval index-one claim.
This makes no assertion about other attractors there.
It is a static full-spectrum counterexample, not a theorem about a
crossing, a Hopf bifurcation or a periodic branch.
\end{remark}

\subsection{Synchronous threshold and the precise local basin statement}
\label{slot:P10-4}

\begin{proposition}[Exact basin threshold on the nonnegative synchronous ray]
\label{prop:P10-ray}
For every $g_c<g<g_f$, the scalar flow on the ray satisfies
\[
 \begin{array}{c|ccc}
 r&(0,r_-)&(r_-,r_+)&(r_+,\infty)\\ \hline
 f(r,g)&-&+&-
 \end{array}
\]
If $0<r(0)<r_-$, then $r(t)\to0$. If $r(0)>r_-$, then
$r(t)\to r_+$; this includes the constant solution $r(0)=r_+$.
The data $r(0)=r_-$ and $r(0)=0$ give their respective constant
solutions. These are exact full-system solutions by synchrony,
even when the upper equilibrium is transversely unstable.
\end{proposition}
\begin{proof}
Factor \eqref{eq:P10-scalar} as
\begin{equation}
 f(r,g)=-\frac{Ar}{1+r}(r-r_-)(r-r_+).
 \label{eq:P10-ray-factor}
\end{equation}
The prefactor is negative for $r>0$, giving the three signs.
Uniqueness prevents crossing an equilibrium. P01 and P02 give global
nonnegative existence and boundedness. More specifically solutions
below $r_-$ decrease and are bounded below by zero; solutions between
the roots increase and are bounded above by $r_+$; those above $r_+$
decrease and are bounded below by it. Each therefore has a limit.
If that limit were not a zero of $f$, continuity would bound $|f|$
away from zero near the limit, contradicting convergence. The monotone
interval and the initial strict inequality select the claimed zero.
No nonsynchronous basin conclusion follows from this argument alone.
\end{proof}

\begin{proposition}[Unstable synchrony and a local separating stable disk]
\label{prop:P10-local-basins}
At a near-fold parameter tuple covered by
Corollary~\ref{cor:P10-near-stability} and
Theorem~\ref{thm:P10-bistability}, let $p=r_-\one$.
The local stable manifold is a smooth embedded $(n-1)$-disk tangent
to $E$. Its germ locally separates a sufficiently small neighborhood
into two sides. The local unstable manifold is exactly a segment of
the synchronous line through $p$. Its two forward branches tend,
respectively, to the origin and to $r_+\one$.

In fact one may choose a neighborhood $V\subset\PosOrth$ of $p$
such that its local stable disk $S$ has
\begin{equation}
 V\setminus S=V_-\mathbin{\dot\cup}V_+,\qquad
 V_-\subset\mathcal A(0),\quad V_+\subset\mathcal A(r_+\one),
 \label{eq:P10-local-partition}
\end{equation}
where the two sets are connected and $\mathcal A(z)$ denotes the
basin of the equilibrium $z$ relative to $\Orth$. Every point of
$S$ tends to $p$. Thus $S$ is the common basin boundary \emph{relative
to this neighborhood}. No assertion identifies the full global basin
boundary or globally continued stable manifold.
\end{proposition}
\begin{proof}
The lower saddle has stable space $E$ and unstable space
$\operatorname{span}\{\one\}$. The stable/unstable manifold theorem
gives the dimensions and tangencies. Tangency and invariance alone
would not establish an arbitrary invariant curve's equality with
an unstable manifold, so check the backward-orbit characterization.
On a sufficiently small synchronous interval around $r_-$ the scalar
field has a simple repelling zero, $f_r(r_-,g)>0$. Every point of
that interval has a negative-time orbit staying in the interval and
converging exponentially to $r_-$ as $t\to-\infty$. Its vector lift
has the same property. Uniqueness of the local unstable manifold,
as the set of such backward trajectories, identifies its germ with
this interval. Both are one-dimensional, so after restriction it is
exactly the stated segment. Proposition~\ref{prop:P10-ray} gives
the two forward limits. The open branches lie on $0<r<r_-$ and
$r_-<r<r_+$ and furnish synchronous heteroclinic trajectories.

Here is a local argument also identifying the two basin sides, rather
than extrapolating a global boundary from codimension. Write
$x=p+s\one+y$, $y\in E$. The stable disk is a graph $s=\phi(y)$
with $\phi(0)=0$, $D\phi(0)=0$. Use the smooth local coordinates
$\xi=s-\phi(y)$, $z=y$, so the stable disk is $\xi=0$ and the
unstable segment is $z=0$. Local invariance of these two manifolds
and smoothness imply
\begin{equation}
 \dot\xi=\xi\bigl(\lambda+\alpha(\xi,z)\bigr),\qquad
 \dot z=Bz+R(\xi,z),
 \quad \lambda=\lambda_0(p)>0,\quad B=J(p)|_E,
 \label{eq:P10-saddle-chart}
\end{equation}
where $\alpha(0,0)=0$, $R(\xi,0)=0$ and $D_zR(0,0)=0$.
P04 normality and the transverse gap give
$\langle z,Bz\rangle\le-\gamma\|z\|^2$ for real $z\in E$,
with some $\gamma>0$. Shrink the coordinate neighborhood so
$|\alpha|<\lambda/2$ and $\|R(\xi,z)\|\le(\gamma/2)\|z\|$.
Hence while in it, $\|z(t)\|\le e^{-\gamma t/2}\|z(0)\|$
and $|\xi(t)|$ grows at least at rate $\lambda/2$ if nonzero.

Choose $\rho>0$ small so $(\pm\rho,0)$ lie on the respective
synchronous unstable branches and the closed coordinate box is in
this neighborhood and in $\PosOrth$. These two exit points belong
to the two sinks' basins. Those basins are relatively open: each sink
has a local attracting neighborhood and continuous dependence pulls
it back along any finite segment of a converging orbit. Thus choose
$\sigma>0$ small enough that both exit caps
$\{(\pm\rho,z):\|z\|\le\sigma\}$ lie in the respective basins.
The image $V$ of $(-\rho,\rho)\times\{\|z\|<\sigma\}$ is an
open neighborhood of $p$. A solution starting with $\xi\ne0$
cannot exit through a $z$ boundary, because its norm decreases,
and reaches the cap with the sign of $\xi$ in finite time. It then
converges to the corresponding sink. If $\xi=0$, it stays on the
stable graph and converges to $p$. The two half-boxes are connected
and map to $V_-,V_+$. Each disk point is approached from both
halves, so the disk is exactly the boundary of each basin relative
to $V$. This proves \eqref{eq:P10-local-partition} with the claimed
local scope. The disk is locally invariant (and forward invariant
in this chosen box), not necessarily invariant under arbitrary
negative-time continuation within a fixed small neighborhood.
\end{proof}

\paragraph{Global scope limit.}
The ray threshold is global only within the invariant synchronous ray.
The preceding full-space statement concerns one sufficiently small
neighborhood of the lower saddle at a fixed near-fold tuple. It neither
identifies the complete global boundary between the sinks nor gives a
partition of the whole orthant. Other equilibria, periodic attractors,
or other invariant attracting sets elsewhere are not excluded. In
particular ``bistability'' here does not assert exactly two global attractors.

\subsection{Connection to P09, degenerate limit, and the P11 boundary}

For $b>0$, as $g\downarrow g_c$ from above,
\begin{equation}
 r_-\longrightarrow0,\qquad r_+\longrightarrow b/A,
 \qquad \frac{dr_-}{dg}=\frac1\Delta\longrightarrow\frac1b.
 \label{eq:P10-transcritical-connection}
\end{equation}
Thus P09's small positive branch on $g>g_c$ is precisely P02/P10's
lower branch, unstable in the uniform direction throughout
$(g_c,g_f)$. It meets the upper branch at $g_f$. No upper stability
all the way to $g_c$ follows; Remark~\ref{rem:P10-counterexample}
shows why this would be false in general.

For fixed $k,c,h>0$, as $b=a-A\downarrow0$,
\begin{equation}
 r_f=\frac b{2A}\to0,\qquad g_f-g_c=\frac{b^2}{4A}\to0,
 \qquad \alpha_{\rm SN},\beta_{\rm SN}\to0.
 \label{eq:P10-degenerate-limit}
\end{equation}
The limiting locus is P09's $a=A$, $g=a+c$, where the center field
is $-Ar^3/(1+r)$, its cubic derivative is $-6A$, and the center
zero remains simple. P09 supplies the codimension-two algebraic/spectral
statement. The P10 nondegeneracy theorem excludes this limit.
There is no two-parameter unfolding or cusp classification here.

The exact facts available for later P11 work are the uniform signs
on the entire open two-root interval, the complete fold spectrum,
and automatic all-mode upper stability sufficiently near the fold.
A subsequent loss of upper stability inside this interval must be
nonuniform. P10 establishes no Hopf existence or frequency, critical
pair simplicity, crossing transversality, Lyapunov coefficient, or
periodic branch. The static counterexample does not begin those tasks.
The frozen direct edge into P10 remains P02/P03/P04/P09, with P01
inherited; recording potential later inputs does not alter the frozen DAG.

\subsection{Checks and independent audit priorities}

B03 is resolved for this DRAFT as stable coexistence, the exact
synchronous threshold, and the explicitly proved local saddle basin
statement; global exclusivity and complete basin geometry remain open.
This fits the frozen S4 supporting role, so no scope amendment is needed.
A separate checker evaluates exact rational identities, original-field
residuals and direct Jacobians, tests the full spectrum at and near the
fold, and searches wider two-root parameters for counterexamples.
It has no theorem status and cannot certify equilibrium completeness
or manifold claims by numerical sampling. All analytical claims above
remain DRAFT pending independent A10 audit.

\paragraph{Priority checks for the independent P10 auditor.}
\begin{enumerate}[leftmargin=*]
\item Recheck the denominator/amplitude notation, square completion,
$g_f-g_c>0$, and the exact branch labels inherited from P02.
\item Derive $G_f-H_f=r_f(k+c/s_f^2)>0$ from the physical fold data;
check the transverse margin for all labels, parity, offset collisions
and noncoprime cases. Distinguish a simple zero from other repetitions.
\item Check the real transverse IFT, its product neighborhood, and
uniqueness forcing $y=0$ for all nearby full equilibria.
\item Recompute $F_g$ and the second directional derivative from the
original field; check $w^{\mathsf T}v=1$ and the half-Taylor convention.
\item Check the parameter-extended nilpotent center block and generalized
center tangent space, the exact invariant center manifold, and the
strictly positive time factor. No full-space decoupling is used.
\item Check the exact uniform signs on the entire two-root interval
and finite-mode transverse persistence on both roots. The index-one
claim is local to the fold, not extended without proof.
\item Check the stable origin, nonempty coexistence interval, smooth
hyperbolic persistence in all five coefficients, openness of
$\mathcal B_{n,q}$, and the sufficient criterion including equality.
\item Check the explicit sufficient interval and the exact unstable
upper-branch counterexample, independently of numerical searches.
\item Check all scalar sign intervals and limiting equilibria; justify
unstable-manifold equality by backward convergence, not tangency alone.
\item Check local stable-disk separation and, especially, the cap-exit
argument for local basin membership. The statement is relative to
$V$ and does not identify the complete global basin boundary.
\item Preserve the P09 degenerate-limit qualification and strict P11
boundary: no Hopf, critical-pair, transversality, nonlinear periodic
or two-parameter unfolding theorem is supplied.
\end{enumerate}
